
Charged-current \numu events must have a reconstructed \mum.
The muon is identified by a minimum ionizing track that
traverses \minerva~\cite{minerva_nim} and travels downstream into the
\minos spectrometer~\cite{Michael:2008bc} where its momentum
and charge are measured.
Muon selection and energy (\Emu) reconstruction are described in Refs.~\cite{nubarprl,nuprl,minerva_nim}.
Requiring a matching track in \minos restricts muon acceptance.
To minimize acceptance differences across the \minerva detector, the analysis
requires neutrino energies above \unit[2]{GeV} and muon angles with respect to the beam (\thetamu) 
less than \unit[17]{$^{\circ}$}.
A \unit[20]{GeV} upper limit on neutrino energy reduces the \numubar background to below 1\%.
After all selection criteria, 5953 events in C, $19\,024$ in Fe, $23\,967$ in Pb, and $189\,168$ in CH are analyzed.

The event vertex is the location of the most upstream energy deposition on the muon track
when only one track is reconstructed; a Kalman filter~\cite{Fruhwirth1987444,Luchsinger1993263} is used to fit the vertex position for events with multiple tracks.
Between 10\% and 20\% of selected events in the different samples have a well-reconstructed multitrack vertex; the remainder are single track
or have a poorly reconstructed vertex position based on the $\chi^2$ of the vertex fit.
Events with vertices in targets 2 through 5 and the fully active tracking volume are considered. 
The target 1 sample has a higher background from interactions upstream of the detector.

Events with a vertex in the active tracking region are divided into three statistically independent CH samples used to form ratios with C, Fe, and Pb.
Events are associated with the C, Fe, or Pb of a nuclear target if the vertex position is between one plane upstream and two planes downstream of that nuclear target module and more than 
 \unit[2.5]{cm} away transversely from seams that join different materials in the target.
In single-track events, the muon track is propagated to the longitudinal center of the nuclear target to estimate the vertex 
position and momentum of the muon.
After all cuts, charged-current event selection efficiency ranges from 24\% in the most upstream 
targets to 50\% in the most downstream. 
The large efficiency variation exists because the upstream region has more inert material and smaller \minos solid angle coverage.   

Energy of the hadronic recoil system \recoilE is determined from the calorimetric sum
of energy deposits not associated with the muon track. 
We consider deposits which occur
between \unit[20]{ns} before and \unit[35]{ns} after the muon to reduce contributions from overlap with other neutrino interactions.
Visible energies are weighted to account for the active fraction of scintillator in different regions of the detector.
The overall calorimetric scale comes from fitting reconstructed \recoilE 
to generated \recoilE for simulated events in the active tracking region~\cite{minerva_nim}.
Using the same procedure, additional calorimetric scales for events in 
targets 2 through 5 are obtained as relative to the tracker;
these are, respectively, 1.11, 1.04, 0.99, and 0.98.  

Kinematic variables \Enu, \Xbj, and \Qsq are obtained 
from reconstructed \Emu, \thetamu, and \recoilE:
$\Enu = \Emu + \recoilE$, $\Qsq = 4 \Enu \Emu {\rm{sin}}^2(\frac{\thetamu}{2})$, 
and $\Xbj = \frac{\Qsq}{2 M_{N}\recoilE}$, where $M_{N}$ is the average 
of proton and neutron masses.
Reconstructed \Enu distributions are corrected for detector smearing 
using iterative Bayesian
unfolding~\cite{D'Agostini:1994zf} with four iterations to produce 
event yields as functions of unfolded \Enu, with generated \Enu values from GENIE.

Reconstructed \Xbj is smeared broadly, especially at high \Xbj where quasielastic processes dominate.
For these events, the hadronic recoil system can be a single nucleon, which is not reconstructed well under a calorimetric assumption.
Such significant smearing would cause large uncertainties in the unfolding procedure.
We therefore report cross section ratios as functions of reconstructed \Xbj\SuppNote{See Supplemental Material\SuppLocation for migration matrices necessary for comparisons of theoretical models to these data}.

Nuclear target samples contain events from adjacent tracking modules due to the loose cut on longitudinal vertex position.
This background, called ``CH contamination,''
ranges from 20\% to 40\% and is roughly proportional to 
the ratio of areal densities of the target to surrounding scintillator.
CH contamination
is estimated by extrapolating event rates measured
in the active tracking region to the nuclear target region.
The tracking and nuclear target regions occupy different areas, and therefore have different acceptance into the \minos detector.
Further, the Fe and Pb targets in the nuclear target region stimulates greater activity in hadronic showers, 
which affects tracking efficiency.
To account for the geometric acceptance difference, we apply a 
correction $w^{t,\A}(\Emu,\thetamu)$,
obtained from a large, single-particle simulated \mum sample.
Here $t=2,3,4,5$ is the target number and \A$=\:$C, Fe, Pb is the target nucleus.
We account for $\recoilE$-dependent efficiency differences using
simulated neutrino events to derive a correction $w^{t,\A}(\recoilE)$.
Differences are largest at low \recoilE.
Acceptance- and efficiency-corrected distributions are scaled such that the integrated number of
events in true and estimated backgrounds are equal according to neutrino event simulation.
Figure~\ref{fig:sample_x} shows the \Xbj distribution of events passing all 
selection criteria in data and simulation; the estimated CH contamination is also shown.

Deviations found in simulated events between the estimated CH contamination by extrapolation and the predicted CH contamination using generator-level information
are not fully covered by statistical uncertainty at the 68\% confidence level in all targets.
Additional systematic uncertainty is applied to ensure coverage at the 68\% level.
CH contamination uncertainties are 1\%--8\% from these systematic deviations and 2\%--5\% from statistics.

\begin{figure}[tp]
\centering
\includegraphics[width=0.825\columnwidth]{SamplePlot_dataMC_x_t02_z82_Nu_v19p1}
\caption{Reconstructed Bjorken $x$ distributions in data and simulation for selected inclusive \numu events in the lead of Target 2. The plot includes CH contamination separately estimated using data and simulated events in the tracker region.  Both simulation distributions are normalized to the data by the number of events passing all event selection criteria.  Events are scaled to a bin size of 0.1.  Events with \Xbj greater than 1.5 are not shown.} 
\label{fig:sample_x}
\end{figure}

Small backgrounds from \numubar ($<0.4\%$) and neutral current ($<0.1\%$) events are estimated using simulation and subtracted.
Transverse smearing within a nuclear target causes roughly 0.5\% of the interactions to be assigned an incorrect target nucleus;
this is also estimated by simulation and subtracted.
A background from upstream neutrino interactions of $6.2\%\pm2.4\%$ exists only in target 2 for one third of the beam exposure, 
because two of the modules upstream of target 1 were not yet instrumented;
affected data are weighted accordingly.

GENIE predicts a sample not dominated by any single process.
Table~\ref{tab:sample_comp_Summary_table} shows the predicted prevalence of processes in bins of reconstructed \Xbj.
%%%
\begingroup
\squeezetable
\begin{table}
\begin{tabular}{cccccccc}
\hline \hline
Reconstructed \Xbj & I    & II   & III  & IV   & V    & Mean Generated \Qsq\\
                   & (\%) & (\%) & (\%) & (\%) & (\%) &  (\GeVcsq)\\
\hline
0.0--0.1 & 11.3    &  42.5    &  5.9   &  19.2  & 15.7  &  0.23 \\
0.1--0.3 & 13.6    &  36.4    &  16.7  &  9.1   & 23.0  &  0.70 \\
0.3--0.7 & 32.7    &  32.8    &  11.8  &  1.4   & 21.1  &  1.00 \\
0.7--0.9 & 55.1    &  25.4    &  4.3   &  0.5   & 14.6  &  0.95 \\
0.9--1.1 & 62.7    &  21.6    &  2.8   &  0.5   & 12.3  &  0.90 \\
1.1--1.5 & 69.6    &  18.1    &  1.9   &  0.4   & 9.9   &  0.82 \\
$>1.5$   & 79.1    &  12.8    &  0.6   &  0.3   & 7.1   &  0.86 \\ 
\hline \hline
\end{tabular}
\caption{
  Average sample composition of selected nuclear target and tracker events in reconstructed \Xbj bins
  based on GENIE simulation of different physics processes, together with the average generated \Qsq.
  Processes are (I) quasielastic, (II) baryon resonance production, (III) deep inelastic scattering at $\Qsq>1$~\GeVcsq and $W>2$~\GeV, (IV) deep inelastic scattering at $\Qsq<1$~\GeVcsq and $W>2$~\GeV, and (V) nonresonant inelastic continuum with $W<2$~\GeV.
} 
\label{tab:sample_comp_Summary_table}
\end{table}
\endgroup
We compare GENIE's prediction for inclusive cross section ratios restricted to $2<\Enu<20~\GeV$ and $\thetamu<17^{\circ}$ to two other models for nuclear modification of structure functions\SuppNote{See Supplemental Material\SuppLocation for a table summarizing the comparison of models of nuclear modification ofinelastic structure functions}.
The Kulagin-Petti microphysical model starts with neutrino-nucleon structure functions and incorporates \A-dependent nuclear effects~\cite{PhysRevD.76.094023,Kulagin2006126}.
The updated Bodek-Yang treatment~\cite{by13} of the model implemented in GENIE~\cite{Bodek:2004pc} includes an A-dependent empirical correction based on charged lepton measurements on the nuclei of interest.
Although nuclear structure functions vary by 20\% among models, ratios of structure functions in Fe or Pb to C differ by $\lesssim$1\%.
